% Generated draft: requires mathematical and source review.
\documentclass[11pt]{article}
\usepackage[margin=25mm]{geometry}
\usepackage{fontspec}
\setmainfont{Latin Modern Roman}
\usepackage{amsmath,amssymb,mathtools,unicode-math}
\setmathfont{Latin Modern Math}
\usepackage{enumitem,needspace,xurl,hyperref}
\setlength{\parindent}{0pt}
\setlength{\parskip}{5pt}
\setlength{\emergencystretch}{3em}
\newcommand{\N}{\mathbb{N}}
\newcommand{\Z}{\mathbb{Z}}
\newcommand{\Q}{\mathbb{Q}}
\newcommand{\R}{\mathbb{R}}
\newcommand{\C}{\mathbb{C}}
\newcommand{\F}{\mathbb{F}}
\newcommand{\A}{\mathbb{A}}
\newcommand{\PP}{\mathbb{P}}
\newcommand{\E}{\mathbb{E}}
\newcommand{\eps}{\varepsilon}
\newcommand{\dd}{\,\mathrm{d}}
\begin{document}
\section*{UnsolvedMath ID: 3900011}
\textbf{Source problem code:} catalogue entry
\subsection{Definitions and mathematical statement}
For a nondegenerate tetrahedron $T=\operatorname{conv}\{v_1,v_2,v_3,v_4\}\subset\R^3$, define its total edge length
\[
 L(T)=\sum_{1\leq i<j\leq4}\|v_i-v_j\|.
\]
For two tetrahedra $T$ and $S$, the ratio in the question is $L(T)/L(S)$, not a volume ratio or an individual-edge ratio. No similarity between the two tetrahedra is stated.

\textbf{Problem.} Determine the largest possible value of $L(T)/L(S)$ when $T$ can ``pass through or fit inside'' $S$ in the intended geometric sense.

The containment part has a usual precise version: $T$ fits inside $S$ if some rigid motion sends $T$ into $S$, so its extremal target is
\[
 \sup\left\{\frac{L(T)}{L(S)}:\exists R\in\operatorname{SO}(3),\ t\in\R^3,
                      \ RT+t\subseteq S\right\}.
\]
Touching is allowed in this displayed version. Whether the source requires strict containment, allows degenerate limiting tetrahedra, or asserts attainment of the supremum is not specified.

The passage part is not defined. A Prince Rupert construction ordinarily requires a specified aperture/tunnel, permissible removed material, and rigid-motion passage convention. The excerpt supplies none, nor does it say whether the containing tetrahedron is treated as a solid, boundary, or framework. Passage therefore cannot be declared equivalent to containment or replaced by an orthogonal-projection condition without source context. It is also unclear whether the two alternatives are to have separate optimum ratios or a single optimum over their union.

The supplied research summary reports computational investigation, claimed characterizations of some Rupert families, and exact ratios for notable cases such as a regular tetrahedron, but says the largest ratio over the whole tetrahedral family was not confirmed from its sources. Those claims concern a Prince Rupert interpretation, whose relation to the stated arbitrary-tetrahedron edge-sum ratio is not specified. They are not independent verification of the target or its status.
\subsection{Short English statement}
Maximize the ratio of total edge lengths of one tetrahedron to another when the first can pass through or fit inside the second. Containment has a standard rigid-motion meaning, but passage is undefined and may be a different optimization problem. Reported Rupert results are not matched to this exact edge-sum objective.
\subsection{Sources}
\begin{itemize}
\item Local UnsolvedMath registry, permanent ID 3900011, source code catalogue entry, ``Prince Rupert ratio for tetrahedra''; attribution field: David Eppstein and cited proposers. Supplied status partially solved, classification PARTIAL-PROGRESS, updated 2026-08-06.
\item Eppstein, \emph{The Geometry Junkyard: Open Problems}, item ``Prince Rupert ratio for tetrahedra,'' supplied index \url{https://ics.uci.edu/~eppstein/junkyard/open.html}. Reported HTML extraction and access: 2026-07-29. The record says item-level later resolution was not checked. Rights: NEEDS\_REVIEW; public access verified but redistribution permission not established.
\item The summary's computational-family and regular-tetrahedron ratio claims have no supplied papers, exact values, theorem statements, or geometric passage conventions.
\end{itemize}
\end{document}
